\par\SourceLeaf{424}%
\noindent\hspace*{1.2em}Der Beweis ist leicht so zu führen. Es sei
\[
AL=t,\quad AL'=t',\quad AL''=t'',\quad AL'''=t''';
\]
man hat dann
\[
\begin{aligned}
\cos LL'&=\cos t\cos t'+\sin t\sin t'\cos A\\
\cos L''L'''&=\cos t''\cos t'''+\sin t''\sin t'''\cos A\\
\cos LL''&=\cos t\cos t''+\sin t\sin t''\cos A\\
\cos L'L'''&=\cos t'\cos t'''+\sin t'\sin t'''\cos A.
\end{aligned}
\]
Also
\[
\begin{aligned}
&\cos LL'\cos L''L'''-\cos LL''\cos L'L'''\\
&=\cos A\{\cos t\cos t'\sin t''\sin t'''+\cos t''\cos t'''\sin t\sin t'\\
&\qquad\qquad-\cos t\cos t''\sin t'\sin t'''-\cos t'\cos t'''\sin t\sin t''\}\\
&=\cos A(\cos t\sin t'''-\cos t'''\sin t)(\cos t'\sin t''-\cos t''\sin t')\\
&=\cos A\sin(t'''-t)\sin(t''-t')\\
&=\cos A\sin LL'''.\sin L'L''.
\end{aligned}
\]

Da jeder der beiden grössten Kreise von $A$ aus in zwei entgegengesetzten Richtungen ausläuft, so bilden sich daselbst zwei einander zu $180^{0}$ ergänzende Winkel; man sieht aber aus unserer Analyse, dass man diejenigen Schenkel wählen muss, die in dem Sinn von $L$ nach $L'''$ und von $L'$ nach $L''$ laufen.

Statt des Winkels $A$ kann man auch die Distanz des Pols des grössten Kreises $LL'''$ vom Pol des grössten Kreises $L'L''$ nehmen; da jedoch jeder grösste Kreis zwei Pole hat, so sieht man, dass man diejenigen verbinden muss, um welche resp. die grössten Kreise von $L$ nach $L'''$ und von $L'$ nach $L''$ in einerlei Sinn laufen.

Die Ausführung des speciellen Falls, wo von den Bögen $LL'''$, $L'L''$ einer oder beide $90^{0}$ gross sind, überlassen wir dem Leser.

6) Noch ein nützlicher Satz ergibt sich aus folgender Analyse. Es seien $L,L',L''$ drei Punkte auf der Kugelfläche und
\[
\begin{aligned}
\cos L(1)&=x,&\quad\cos L(2)&=y,&\quad\cos L(3)&=z,\\
\cos L'(1)&=x',&\cos L'(2)&=y',&\cos L'(3)&=z',\\
\cos L''(1)&=x'',&\cos L''(2)&=y'',&\cos L''(3)&=z''.
\end{aligned}
\]
